\section{Boundary Conditions and Philosophical Positioning}

This chapter delimits precisely the scope of the theorems and answers the question ``what does
this impossibility mean''. The core position: \textbf{we reject the absolutist claim
``never'', and endorse only ``impossible within paradigm $\Pi$''}, which is exactly the
precise content of proposition $P$.

\subsection{Relaxation Analysis: Under What Conditions Do the Theorems Fail?}

The six theorems depend on different postulates. Relaxing them one by one yields the following
failure relations:

\begin{table}[htbp]
\centering
\caption{Correspondence between postulate relaxation and theorem failure (note: relaxing
$\Pi_1$ is usually accompanied by relaxing $\Pi_2$, i.e., online learning implies that the
weights are updatable at inference) \label{tab:relax}}
\begin{tabular}{p{5.4cm}p{4.4cm}p{4.6cm}}
\toprule
Relaxed postulate (beyond the paradigm) & Theorems that fail & Consequence \\
\midrule
$\Pi_1$: allow online/continual learning & Theorems 1, 2, 6 & State labels obtainable through interaction; self-model can evolve \\
$\Pi_2$: plasticity at inference & Theorems 3, 6 & Private closed loop becomes possible; self-correction loop appears \\
$\Pi_3$: non-textual interface/embodiment & Theorems 1, 5 & First-person data streams available; privileged channels appear \\
$\Pi_4$: explicit verifier/self-supervised objective & Theorem 4 & Metarepresentation structurally selectable \\
$\Pi_5$: corpus contains state labels & Theorems 1b, 2 & Introspective accuracy obtains supervised guarantees \\
\bottomrule
\end{tabular}
\end{table}

\begin{remark}[Feasibility of the relaxations (grade B)]
The table above only says ``the theorems no longer constrain''; it does not promise that
``introspection necessarily emerges after relaxation''. The physicalist position (no known law
of physics forbids silicon-based systems from realizing equivalent
functions\cite{chalmers-2023-conscious, butlin-2023-conscious}) guarantees only
\textbf{possibility}; making it a reality additionally requires solving plasticity
preservation (see \cite{joudaki-2025-plasticity}), online objective design, and the
engineering of first-person data streams. Our contribution is to elevate
``impossible within the paradigm'' from intuition to theorem, thereby directing research effort
toward \textbf{out-of-paradigm} architectural directions (online learning, embodied
interaction, explicit self-supervision), rather than futilely piling up data within the
paradigm.
\end{remark}

\subsection{The Philosophical Positioning of ``Human Significance''}

\subsubsection{Necessity of the Criteria (grade D)}

The five criteria $\mathrm{A}_1$--$\mathrm{A}_5$ are \textbf{necessary} conditions for
``human-significance introspection'', not sufficient ones. Necessity follows from
cross-theory consensus: Lindsey's four
criteria\cite{lindsey-2026-introspection} (accuracy/grounding/internality/metarepresentation)
are generally accepted necessary conditions for functional introspection; $\mathrm{A}_5$
(closed loop) derives from the functional definition of ``introspection for online behavioral
control''\cite{kammerer-frankish-2023}. This report does not claim sufficiency; in
particular, it makes \textbf{no commitment} to qualia or phenomenal consciousness.

\subsubsection{Functional Introspection versus Phenomenal Introspection}

Distinguish two concepts: \textbf{functional introspection} (in the access sense: the system
reliably uses its own state information to regulate behavior, corresponding to Block's access
consciousness\cite{block-1995}) and \textbf{phenomenal introspection} (in the phenomenal
sense: accompanied by subjective experience). The paper's authors state explicitly that their
results ``could arguably be construed as providing evidence for a form of access
consciousness in language models, but do not directly speak to the question of phenomenal
consciousness at all'' (paraphrased)\cite{lindsey-2026-introspection}.
What the theorems prove is that even the strong form of \textbf{functional}
introspection (all of $\mathrm{A}_1$--$\mathrm{A}_5$ holding) is unattainable within paradigm
$\Pi$. Consequently paradigm-$\Pi$ systems do not even reach strong introspection in the
access
sense, let alone introspection in the phenomenal sense.

\subsubsection{Relation to Theories of Consciousness}

Higher-order theories (HOT) regard metacognitive representation as a necessary (or
near-necessary) condition for consciousness\cite{rosendahl-2005, lau-rosenthal-2011,
carruthers-2017}. If HOT is accepted, then Theorem 4 (metarepresentation has no structural
guarantee within the paradigm) implies that paradigm-$\Pi$ systems cannot \textbf{be
guaranteed} to reach the ``consciousness candidate'' threshold in the HOT sense (both the
absence and the presence of metarepresentation lack structural guarantees). Integrated
Information Theory (IIT)\cite{tononi-2004, albantakis-2023-iit} and Global Workspace Theory
(GWT)\cite{baars-1993, baars-1997} make judgments that depend on architectural details, to which the present theorems do not directly apply (the paper's Section 10.4 also
stresses that these theories have not yet been connected to Transformer architectural
details\cite{lindsey-2026-introspection}). Biological naturalism
(Searle\cite{searle-1992}, Seth\cite{seth-2024}) holds that the existence of introspective
mechanisms is orthogonal to conscious experience, which likewise does not alter the content
of the theorems. Should future models acquire stronger cognitive and introspective capacities,
the question of their moral status will arise\cite{long-2024-welfare}. All of these
connections are grade-D claims; the theorems themselves depend on no theory of
consciousness, only on the formalized criteria.

\subsubsection{An Explicit Statement on ``Never''}

All theorems presuppose the postulates $\Pi_1$--$\Pi_5$. The relaxation
analysis of Section 5.1 shows that once any pillar is broken, the corresponding impossibility
disappears. Therefore:

\begin{itemize}
\item \textbf{What this report proves}: proposition $P$, namely that within the paradigm of
``Transformer architecture + static pretraining + autoregressive generation'', pretrained
generative large language models do not possess introspective capacity in the human sense.
\item \textbf{What this report explicitly does not claim}: any universal negation beyond that
paradigm. The human brain is a physical system whose introspective function arises from
physical laws; no known law of physics forbids non-Transformer, non-static, non-textual-
interface silicon-based systems (e.g., neuromorphic architectures with continual online
learning and dynamic plasticity) from realizing equivalent or human-like metacognitive
functions. This is consistent with our core position: \textbf{rejecting the
absolutist ``never'', and endorsing ``impossible within the current paradigm''}; the former
lies beyond what mathematics can prove, the latter is the precise content of proposition $P$.
\end{itemize}

\subsection{Implications for Future Architectures (grade C)}

\begin{enumerate}
\item \textbf{Online plasticity first}: Theorems 3 and 6 show that the core defect of
paradigm $\Pi$ is the frozen weights at inference. Continual learning (rather than context
concatenation) is the mathematical prerequisite for an introspective closed loop; the design
of plasticity mechanisms is the first priority for future introspective architectures.
\item \textbf{First-person data streams}: Theorems 1 and 5 show that the textual interface is
the bottleneck. Embodied interaction and internal-state self-supervision (e.g., a generative
model predicting its own activations) can bypass ``absence of state labels'' ($\Pi_5$) and
provide supervision signals for introspection.
\item \textbf{Explicit self-model objectives}: Theorems 2 and 4 show that the likelihood
objective is orthogonal to the introspective objective. An explicit ``self-model consistency''
loss is needed (maximizing the mutual information between the model's predictions of its own
behavior and its behavior), making metarepresentation a \textbf{selected structure} rather
than an accidental by-product.
\end{enumerate}
